\documentclass[11pt]{article}

\usepackage[a4paper,margin=1in]{geometry}
\usepackage{amsmath,amssymb}

\title{\textbf{CSE400 – Fundamentals of Probability in Computing}\\
\vspace{0.3cm}
\large Lecture 6: Discrete Random Variables, Expectation, and Problem Solving}
\date{}
\author{}

\begin{document}
\maketitle

\noindent
\textbf{Source:} Lecture 6 Slides

\vspace{0.3cm}

\section{Random Variables}

\subsection{Motivation and Concept}

A \textbf{random variable} is introduced to assign numerical values to outcomes of a random experiment.

\subsubsection*{Definition}
A random variable $X$ on a sample space $\Omega$ is a function
\[
X : \Omega \rightarrow \mathbb{R}
\]
which assigns to each sample point $\omega \in \Omega$ a real number $X(\omega)$.

\subsubsection*{Reasoning}
\begin{itemize}
    \item Outcomes in a sample space may be non-numeric.
    \item To analyze randomness mathematically, outcomes are mapped to real numbers.
    \item This mapping allows the use of algebra, summation, and probability laws.
\end{itemize}

\subsection{Discrete Random Variables}

Until further notice, we restrict attention to \textbf{discrete random variables}.

A random variable is \textbf{discrete} if:
\begin{itemize}
    \item It takes values from a finite or countably infinite set.
    \item Even though $X$ maps into $\mathbb{R}$, the actual values $\{X(\omega)\}$ form a discrete subset of $\mathbb{R}$.
\end{itemize}

\subsubsection*{Reasoning}
\begin{itemize}
    \item Discreteness allows probabilities to be assigned to individual values.
    \item This leads directly to the definition of a Probability Mass Function (PMF).
\end{itemize}

\section{Distribution of a Discrete Random Variable}

The distribution of a discrete random variable can be visualized using a \textbf{bar diagram}.
\begin{itemize}
    \item The $x$-axis represents possible values of the random variable.
    \item The height of each bar at value $a$ equals
    \[
    \Pr[X = a].
    \]
\end{itemize}

\subsubsection*{Reasoning}
\begin{itemize}
    \item Each bar corresponds to an event in the sample space.
    \item The height represents the probability of that event.
    \item This visualization reflects how probability mass is distributed.
\end{itemize}

\section{Example 1: Tossing Three Fair Coins}

\subsection{Experiment Description}

An experiment consists of tossing \textbf{three fair coins}.  
Let $Y$ denote the \textbf{number of heads} that appear.

\subsection{Possible Values of the Random Variable}

The number of heads possible is
\[
Y \in \{0,1,2,3\}.
\]

\subsubsection*{Reasoning}
\begin{itemize}
    \item Zero heads to three heads are the only possible outcomes.
    \item Therefore, the range of $Y$ is finite and discrete.
\end{itemize}

\subsection{Computing the Probabilities}

Each coin toss is fair and independent, so all $2^3=8$ outcomes are equally likely:
\begin{align*}
\Pr(Y=0) &= \Pr(t,t,t) = \frac{1}{8}, \\
\Pr(Y=1) &= \Pr((t,t,h),(t,h,t),(h,t,t)) = \frac{3}{8}, \\
\Pr(Y=2) &= \Pr((t,h,h),(h,t,h),(h,h,t)) = \frac{3}{8}, \\
\Pr(Y=3) &= \Pr(h,h,h) = \frac{1}{8}.
\end{align*}

\subsection{Validity of the Distribution}

Since $Y$ must take one and only one of the values $0,1,2,3$,
\[
\Pr\!\left(\bigcup_{i=0}^{3}\{Y=i\}\right)
= \sum_{i=0}^{3}\Pr(Y=i) = 1.
\]

\subsubsection*{Reasoning}
\begin{itemize}
    \item The events $\{Y=i\}$ are mutually exclusive.
    \item Together, they cover the entire sample space.
    \item Their probabilities must therefore sum to $1$.
\end{itemize}

\section{Probability Mass Function (PMF)}

\subsection{Definition}

Let $X$ be a discrete random variable with range
\[
R_X=\{x_1,x_2,x_3,\ldots\}.
\]
The function
\[
p_X(x_k)=\Pr(X=x_k), \quad k=1,2,3,\ldots
\]
is called the \textbf{Probability Mass Function (PMF)} of $X$.

\subsection{Fundamental Property of PMF}

Since $X$ must assume one of its possible values,
\[
\sum_{k=1}^{\infty} p_X(x_k)=1.
\]

\subsubsection*{Reasoning}
\begin{itemize}
    \item Each value corresponds to a distinct event.
    \item These events form a partition of the sample space.
    \item Total probability must equal $1$.
\end{itemize}

\section{PMF Example with Parameter $\lambda$}

\subsection{Given Information}

The PMF is
\[
p(i)=c\frac{\lambda^i}{i!}, \quad i=0,1,2,\ldots
\]
where $\lambda>0$.

\subsection{Determining the Constant $c$}

Using normalization,
\[
\sum_{i=0}^{\infty} p(i)=1
\quad\Rightarrow\quad
c\sum_{i=0}^{\infty}\frac{\lambda^i}{i!}=1.
\]
Since
\[
\sum_{i=0}^{\infty}\frac{\lambda^i}{i!}=e^{\lambda},
\]
we get
\[
c=e^{-\lambda}.
\]

\subsection{Computing $\Pr[X=0]$}

\[
\Pr[X=0]=e^{-\lambda}.
\]

\subsection{Computing $\Pr[X>2]$}

\[
\Pr(X>2)=1-\Pr(X\le2)
=1-\left(e^{-\lambda}+\lambda e^{-\lambda}+\frac{\lambda^2}{2}e^{-\lambda}\right).
\]

\section{Bayes' Theorem (Recap)}

\subsection{Conditional Probability Identity}

\[
\Pr(A|B)=\Pr(B|A)\Pr(A).
\]

\subsection{Bayes' Formula}

For a partition $\{B_1,B_2,\ldots,B_n\}$,
\[
\Pr(B_i|A)=
\frac{\Pr(A|B_i)\Pr(B_i)}
{\sum_{j=1}^{n}\Pr(A|B_j)\Pr(B_j)}.
\]

\subsection{Interpretation}

\begin{itemize}
    \item $\Pr(B_i)$: Prior probability
    \item $\Pr(B_i|A)$: Posterior probability
\end{itemize}

\section{Bayes' Theorem Example: Auditorium}

\subsection{Problem Setup}

\begin{itemize}
    \item Auditorium has $30$ rows
    \item Row $1$ has $11$ seats, row $30$ has $40$ seats
    \item A row is selected uniformly at random
    \item A seat is selected uniformly from that row
\end{itemize}

\subsection{Computing $\Pr(S_{15}\mid R_{20})$}

Row $20$ has $30$ seats:
\[
\Pr(S_{15}\mid R_{20})=\frac{1}{30}.
\]

\subsection{Computing $\Pr(S_{15})$}

Seat $15$ exists in rows $5$ through $30$:
\[
\Pr(S_{15})
=\sum_{k=5}^{30}\Pr(S_{15}\mid R_k)\Pr(R_k)
=\sum_{k=5}^{30}\frac{1}{k+10}\cdot\frac{1}{30}
=0.0342.
\]

\subsection{Computing $\Pr(R_{20}\mid S_{15})$}

\[
\Pr(R_{20}\mid S_{15})
=\frac{\Pr(S_{15}\mid R_{20})\Pr(R_{20})}{\Pr(S_{15})}
=\frac{\frac{1}{30}\cdot\frac{1}{30}}{0.0342}
\approx0.0325.
\]

\section*{End of Lecture 6 Scribe}

\end{document}
